\NewContentPage
\ReportSubsection{6.3 Bug Findings and Remaining Risks}
The 24 failed examples relate to the nine findings below. All carry \Path{@bug},
and 16 also carry \Path{@policy}. Each finding describes the expected and actual
behavior, its impact, and the affected stories. The full reproduction steps are
recorded in \Path{BUG_REPORT.md}.

\begin{ReportTable}{@{}P{.09\linewidth}P{.11\linewidth}YP{.06\linewidth}@{}}
\toprule
\textbf{ID} & \textbf{Story} & \textbf{Observed outcome and consequence} & \textbf{Fails}\\
\midrule
BR01 & 15 & Missing/deleted project task reads return unrelated work; students can see tasks outside the selected project. & 4\\
BR03 & 13, 14 & Category membership appears only on the writing side; reverse browsing can hide filed work. Reciprocity is a policy. & 6\\
BR04 & 13 & A deleted category remains readable through a task link after direct GET returns 404. & 2\\
BR06 & 11 & Plain XML IDs return 404 while equivalent CDATA and JSON strings link successfully. & 2\\
BR19 & 11 & Numeric JSON IDs return 404 while string IDs work. Numeric coercion is a policy. & 2\\
BR32 & 15 & Title-only project PUT removes task links. Preserving omitted links is a policy. & 2\\
B-N1 & 12 & Title-only task PUT removes its project link; a rename can lose filing. Preservation is a policy. & 2\\
B-N2 & 5 & Invalid done-status filters return 200 with an empty list. Rejecting them is a policy. & 2\\
B-N3 & 5 & Unknown query fields are ignored and return unfiltered work. Rejecting them is a policy. & 2\\
\bottomrule
\end{ReportTable}

\ReportSubsection{6.4 New Findings and Their Consequences}
\RunIn{B-N1: Title-only task replacement removes its project link}
We created a project with two tasks and sent a PUT containing only a new title
to one task. Reading the project's tasks and the task's reverse \Path{tasksof}
view showed that the selected link had been removed, while the other task
remained. Our workflow expects a wording change to preserve project membership,
but the API does not clearly promise this for PUT. Sending the title change
through POST preserved the link.

\RunIn{B-N2: Invalid done-status filter is accepted}
With unfinished and completed todos present, we compared valid boolean filters
with \Path{doneStatus=maybe} and \Path{doneStatus=not-a-boolean}. The invalid
queries returned HTTP 200 and an empty list. Our chosen rule expects HTTP 400,
because an invalid filter should be distinguishable from having no work left.
The API rejects invalid booleans when creating a todo, but that does not establish
the same validation rule for queries.

\RunIn{B-N3: Unknown todo filter is ignored}
We captured the todo list, then queried \Path{colour=red} and \Path{doneStats=red}.
Both returned the same unfiltered records. A valid title filter returned one
record. Our workflow expects HTTP 400 for unknown fields, although the
documentation does not specify that rule. Ignoring a misspelled field can make
a student think a filter was applied when it was not.

\RunIn{What needs attention}
We would prioritize unrelated project data and references to deleted categories,
because both can cause the student to act on the wrong records. Additional
checks reproduced the XML difference, a deleted category in a project
view, and the unknown-filter results~[\ref{ref:probes}]. BR17 records supported
inline creation in stories 11 and 13.
